\frfilename{mt25.tex}
\versiondate{31.5.03}

\def\chaptername{Ukuran produk}
\newchapter{25}
\def\chaptername{Ukuran produk}

Sekarang saya sampai pada bab lain tentang teori ukur `murni', yang membahas
suatu konstruksi fundamental—atau, jika Anda lebih suka memandangnya demikian,
dua konstruksi, karena persoalan dalam membentuk produk dari dua ruang ukur
sebarang (\S251) cukup berbeda dari persoalan yang muncul pada produk dari
sembarang banyak ruang probabilitas ($\S$254). Pekerjaan ini akan menguji teknik
kita sampai batasnya: teorema-teorema fundamental yang hendak kita capai
memang merupakan sasaran yang alamiah, tetapi pembuktiannya panjang dan ada
banyak jebakan di samping jalan yang benar.

\allowmorestretch{500}{
Gagasan utamanya adalah `integrasi berulang'. Anda mungkin sudah pernah
melihat rumus-rumus berbentuk\penalty-100
`$\iint f(x,y)dxdy$' yang digunakan untuk menghitung integral fungsi dua
variabel real di suatu daerah pada bidang. Salah satu teknik dasar kalkulus
lanjut adalah membalik urutan integrasi; misalnya, kita mengharapkan
$\int_0^1(\int_y^1f(x,y)dx)dy$ sama dengan
$\int_0^1(\int_0^xf(x,y)dy)dx$. Dalam pengembangan materi ini, kita sudah
memiliki perhitungan ketiga untuk dibandingkan dengan kedua perhitungan itu:
$\int_Df$, dengan $D=\{(x,y):0\le y\le x\le 1\}$ dan integral diambil
terhadap ukuran Lebesgue pada bidang. Dua bagian pertama bab ini dikhususkan
untuk menganalisis hubungan antara ukuran Lebesgue satu dimensi dan dua
dimensi yang membuat operasi-operasi tersebut sah—setidaknya dalam beberapa
keadaan; sebagian pekerjaan harus digunakan untuk menguraikan dengan cermat
syarat-syarat tepat yang harus dikenakan pada $f$ dan $D$ agar kita aman.
}%end of allowmorestretch

Integrasi berulang, dalam satu bentuk atau bentuk lain, muncul di seluruh
teori ukur; karena itu cepat atau lambat kita perlu mengembangkan ungkapan
yang paling umum untuk gagasan ini. Metode standarnya adalah melalui teori
produk ruang-ruang ukur umum. Diberikan ruang ukur $(X,\Sigma,\mu)$ dan
$(Y,\Tau,\nu)$, tujuannya adalah mencari suatu ukuran $\lambda$ pada
$X\times Y$ yang setidaknya memberikan ukuran yang tepat
$\mu E\cdot\nu F$ kepada sebuah `persegi panjang' $E\times F$ jika
$E\in\Sigma$ dan $F\in\Tau$. Ternyata sudah ada kesulitan dalam menentukan
apa yang dimaksud dengan `ukuran produk', dan untuk mengerjakannya dengan
benar saya perlu, bahkan pada tahap ini, menguraikan dua konstruksi yang
berkaitan tetapi dapat dibedakan. Konstruksi-konstruksi tersebut dan sifat
elementernya mengisi seluruh \S251. Dalam \S252 saya beralih ke integrasi
atas produk, dengan teorema Fubini dan Tonelli yang menghubungkan
$\int fd\lambda$ dengan $\iint f(x,y)\mu(dx)\nu(dy)$. Karena konstruksi
$\lambda$ simetris terhadap kedua faktornya, secara otomatis tersedia teorema
yang menghubungkan $\iint f(x,y)\mu(dx)\nu(dy)$ dengan
$\iint f(x,y)\nu(dy)\mu(dx)$. \S253 membahas ruang $L^1(\lambda)$ dan
hubungannya dengan $L^1(\mu)$ serta $L^1(\nu)$.

Untuk ruang ukur umum, ada hambatan dalam membentuk produk tak hingga; sebagai
permulaan, jika $\sequencen{(X_n,\mu_n)}$ adalah barisan ruang ukur, maka
suatu ukuran produk $\lambda$ pada $X=\prod_{n\in\Bbb N}X_n$ seharusnya
memenuhi $\lambda X=\prod_{n=0}^{\infty}\mu_nX_n$, dan tidak ada jaminan
bahwa produk tersebut konvergen atau berperilaku baik ketika ia konvergen.
Namun, untuk ruang probabilitas, ketika $\mu_nX_n=1$ untuk setiap $n$,
masalah ini setidaknya lenyap. Produk dari keluarga ruang probabilitas apa
pun dapat didefinisikan; inilah pokok bahasan \S254.

Saya mengakhiri bab ini dengan tiga bagian yang menjadi persiapan untuk Bab
27 dan 28, tetapi juga penting pada dirinya sendiri sebagai penyelidikan
tentang cara struktur grup $\BbbR^r$ berinteraksi dengan ukuran Lebesgue dan
ukuran-ukuran lainnya. \S255 membahas `konvolusi' $f*g$ dari dua fungsi,
dengan $(f*g)(x)=\int f(y)g(x-y)dy$ (integralnya terhadap ukuran Lebesgue).
Dalam \S257 saya menunjukkan bahwa beberapa gagasan yang sama, setelah
ditransformasikan dengan tepat, dapat digunakan untuk menguraikan konvolusi
$\nu_1*\nu_2$ dari dua ukuran pada $\BbbR^r$; sebagai persiapan untuk ini
saya menyertakan bagian tentang ukuran Radon pada $\BbbR^r$ (\S256).

\discrpage
